\section{Introduction}
\label{sec:introduction}

Multidimensional coherent spectroscopy provides direct access to
correlations between excitation and detection frequencies and has
become an important tool for resolving electronic couplings,
coherences, relaxation pathways, and many-body interactions in
molecular and condensed-phase systems
\cite{
Abramavicius2009CoherentMultidimensional,
Oliver2018MultidimensionalReview,
Hybl1998TwoDimensional}.
Interpreting these spectra generally requires comparison with a
microscopic dynamical model, making efficient numerical calculation of
nonlinear optical response functions an important part of the
spectroscopic analysis.

Theoretical treatments are often simplified by assuming impulsive
optical interactions and applying the rotating-wave approximation
(RWA).
These approximations provide a compact pathway picture and are useful
over broad regions of parameter space, but finite pulse duration can
introduce effects that are absent from the impulsive description.
Finite-pulse treatments have demonstrated modifications of spectral
amplitudes, phases, line shapes, and observed dynamical signatures,
with the importance of these effects depending on pulse bandwidth,
resonance conditions, and the molecular dynamics occurring during the
interaction
\cite{
Smallwood2017FinitePulse,
Perlik2017FinitePulse,
DoGelinTan2017FinitePulse,
Leng2017FinitePulseWidth}.
For sufficiently broadband fields, the spectral suppression of
counter-rotating light--matter processes can become less effective,
allowing non-RWA contributions to become appreciable
\cite{Smallwood2017FinitePulse}.
Consequently, the validity of the impulsive and rotating-wave
approximations must ultimately be assessed relative to the molecular and
optical timescales of the system under study.

A variety of analytical and numerical strategies have been developed
for treating nonlinear spectroscopy outside the simplest impulsive
picture.
Mixed time--frequency formulations and finite-pulse response
constructions provide alternatives to direct multidimensional
time-domain propagation
\cite{
GelinDomcke2016AlternativeView,
DoGelinTan2017FinitePulse}.
Frequency-domain Green-function approaches have also been used to
describe nonlinear multi-exciton signals
\cite{Roslyak2013Quasiparticle}.
For open quantum systems, the UF$^2$ method evaluates nonlinear signals
through Fourier convolution in the eigenbasis of the Liouvillian and
can efficiently incorporate arbitrary sampled pulse envelopes
\cite{RoseKrich2021UF2}.
Automated diagram generation further provides a systematic treatment
of the additional temporal orderings that arise when pulses overlap
\cite{RoseKrich2021DiagramGenerator}.
Open-source packages such as QuDPy and QDT provide complementary
frameworks for constructing and evaluating multidimensional optical
response functions
\cite{shah2023qudpy,Kenneweg2024QDT}.

The present work addresses a different but complementary computational
organization of the problem.
We consider whether the coherence-frequency coordinates can be
evaluated directly, without first constructing a two-dimensional
time-domain interferogram, while retaining molecular evolution during
finite Gaussian pulse envelopes and, when required, the complete
light--matter interaction.
The treatment remains perturbative in the optical field; strong-field
and nonperturbative two-dimensional spectroscopy constitute a distinct
regime that has been considered elsewhere
\cite{
AndaCole2021BeyondPerturbative,
BressanVanThor2021IntenseFields}.

Here we introduce PULSUS, a Python implementation of a direct
frequency-domain Liouville-space formulation for multidimensional
spectroscopy of open quantum systems.
The coherence-frequency intervals are represented through Liouvillian
resolvents, while finite Gaussian pulses are incorporated through
operator-valued functions of the field-free Liouvillian.
This construction permits the response to be evaluated directly on
either Cartesian frequency grids or arbitrary sets of paired spectral
coordinates.
The complete interaction superoperator may be retained, while the RWA,
short-pulse, and impulsive limits arise as controlled reductions of the
same formulation.

We validate the finite-pulse implementation against the analytical
Gaussian-pulse benchmark of Smallwood \textit{et al.} and verify
agreement between Cartesian-grid and arbitrary-coordinate evaluations.
We then demonstrate targeted spectral sampling in which the evaluated
frequency coordinates are concentrated toward regions containing
significant spectral weight.
Finally, using a coupled electronic dimer, we construct a hierarchy of
four models that separately isolates the effects of the RWA,
molecular evolution during the pulse envelopes, and finite spectral
bandwidth.
The resulting regime map shows that the total error of the
impulsive-RWA approximation is nonmonotonic in pulse duration and is
smallest in an intermediate regime.
For ultrashort pulses, the dominant
nonrephasing beyond-RWA correction is traced to a coupling-enabled
molecular-interaction sequence
$\mathbf{m}=(+,+,-)$ that is removed by the conventional RWA.
Together, these results establish PULSUS as a direct
frequency-domain computational tool and as a framework for diagnosing
which approximations control a given multidimensional spectrum.